% Correlated label fibers for a prime-field quadratic bank.
\subsection{A translated grid with a polynomial source gap}
\label{subsec:translated-grid-polynomial-gap}

Independent fresh values give a logarithmic source gap in the quadratic
bank.  A translated integer grid instead makes many labels share a
polynomial number of fresh coordinates.  The construction below keeps
message dimension three and uses only prime-field coordinates and values.
Its exceptional fraction tends to zero; the assertion is a larger
absolute gap with a superlinear number of singleton labels.

\begin{theorem}[Translated-grid quadratic bank]
\label{thm:translated-grid-polynomial-gap}
There is a sequence of primes $p$ and length-$n$, dimension-three
Reed--Solomon codes over $\mathbb F_p$, with two received words $F,G$
and integers $A<T=A+d$, such that
\[
 \operatorname{agr}_3(F)=\operatorname{agr}_3(G)
   =\operatorname{CA}_3(F,G)=A,
 \qquad
 n a_1(3/n)<A<T<\sqrt{2n}.
\]
At least $B$ distinct labels $\lambda\in\mathbb F_p\setminus\{0,1\}$
have a singleton list at agreement threshold $T$ for
$(1-\lambda)F+\lambda G$, where
\[
 \begin{split}
 d&=\Theta\!\left(\frac{n^{1/6}}{\log n}\right),\qquad
 B=\Omega(n^{4/3}\log n),\\
 p&=\Theta\bigl(n^{4/3}(\log n)^4\bigr),\qquad
 \frac{T-A}{n}
   =\Theta\!\left(\frac{n^{-5/6}}{\log n}\right).
 \end{split}
\]
The successful witness need not have exactly $T$ agreements.
At the relaxed agreement threshold $A$, every member of the affine
line has a list of size at most $L+1$, where $L=\Theta(\sqrt n)$.
Thus this is a vanishing-rate construction with polynomial absolute
source separation; it does not assert a positive constant fraction of
all prime-field labels.
\end{theorem}
\begin{proof}
We first construct the bank and then the correlated fresh coordinates.

\paragraph{A multiplicative Sidon bank.}
Let an integer $H$ tend to infinity.  The prime number theorem in the
fixed progression $1\pmod4$ supplies, for all sufficiently large $H$,
a prime
\[
 2H^4<p<4H^4,\qquad p\equiv1\pmod4.
\]
There are $\Theta(H/\log H)$ primes in $[H/2,H]$.
One of the four quartic-character classes modulo $p$ contains at
least a quarter of them.  Choose a power of two $q_0$ of order
$\sqrt{H/\log H}$ such that this class contains at least
$2v$ primes, where $v=q_0^2+q_0+1$.  Assign disjoint sets of $v$
such primes to the points and lines of the projective plane over
$\mathbb F_{q_0}$.  Join their assigned primes $a,b$ whenever the
point lies on the line.  This bipartite graph has no four-cycle and
has
\[
 L=(q_0+1)(q_0^2+q_0+1)
   =\Theta\bigl((H/\log H)^{3/2}\bigr)
\]
edges.  For each edge choose a square $\theta\in\mathbb F_p^*$ with
$\theta^2=a/b$.  Such a choice is possible because $a/b$ is a
fourth power.  Distinct edges give distinct $\theta$.

These parameters form a multiplicative Sidon set, allowing repeated
pairs.  Indeed, an equality of two products of parameters, after
squaring and clearing denominators, gives a congruence between
positive integers at most $H^4$.  It is an integer equality since
$p>H^4$.  Unique factorization, together with the disjoint numerator
and denominator prime sets, forces the same edge multiset unless the
four edges form a four-cycle.  The latter is excluded.

Take the bank
\[
 P_\theta(X)=\theta+X^2/\theta.
\]
Two distinct bank words meet at precisely the two points
$x^2=\theta\phi$, both in $\mathbb F_p^*$ because the parameters
are squares.  Sidonicity makes these pairs of coordinates disjoint.
The resulting core has $N_0=L(L-1)$ points.  Set $g=0$ on it and
set $f$ to the corresponding common value $\theta+\phi$.
Every bank word has exactly
\[
 A=2L-2
\]
core matches.  Any other quadratic has at most $L$ core matches:
each such match is counted against two bank words, and its difference
with each bank word has at most two roots.

\paragraph{Retaining a translated grid.}
Put $M=\lfloor L/4\rfloor$, and choose $u_0,v_0$ uniformly and
independently in $\mathbb F_p$.  For each $1\le u,v\le M$, write
$U=u_0+u$, $V=v_0+v$.  Keep a grid pair only if $U,V\ne0$ and
$V/U$ is a nonzero square.  Choose one square root $x$ by a fixed
rule.  Retain one representative of every repeated ratio, and delete
ratios belonging to core coordinates.  The resulting fresh coordinates
are distinct; each keeps its associated pair $(U,V)$.

Let $X$ denote their number.  A fixed grid pair is initially eligible
with probability $(p-1)^2/(2p^2)$.  Equality of ratios from different
grid pairs is a nontrivial affine equation in $u_0,v_0$, of probability
at most $1/p$.  Hitting a specified core squared-coordinate also has
probability at most $1/p$.  Therefore
\[
 \mathbb E X\ge
 M^2\frac{(p-1)^2}{2p^2}
 -\frac{\binom{M^2}{2}}p
 -\frac{M^2\binom L2}p
 =\left(\frac12-o(1)\right)M^2.
\]
Since $0\le X\le M^2$, the event $X\ge M^2/4$ has probability
at least $1/3-o(1)$.

At a retained coordinate define
\[
 g(x)=1/U,\qquad f(x)=c_0/U,
\]
where the common shift $c_0$ will be chosen later.  If the bank
parameter satisfies $\theta^2=a/b$, its possible labels obey
\begin{equation}\label{eq:translated-grid-label}
 \lambda+c_0
   =\theta u_0+v_0/\theta+\frac{au+bv}{b\theta}.
\end{equation}
For every bank these raw labels lie in an affine image of the same
integer interval of size $K=3HM<p$, because $1\le au+bv\le2HM$.
An integer value of $au+bv$ occurs in at most
\[
 1+\frac{2M}{H}\le\frac{3M}{H}
\]
grid pairs, for large $H$.  Here $a,b$ are distinct primes at least
$H/2$, so solutions differ by integer multiples of $(b,-a)$.
The bound survives all deletions.

Set
\[
 d=\left\lfloor\frac{M}{64H}\right\rfloor,
\]
which tends to infinity.  If $X\ge M^2/4$, labels having fewer than
$d$ retained matches account for at most $K(d-1)\le3M^2/64$ points.
Dividing the remaining points by the preceding fiber upper bound
shows that every bank has at least $HM/16$ labels with at least
$d$ fresh matches.  The same retained-point count works for all
banks; no independence between their fibers is required.

\paragraph{Distinct labels and far endpoints.}
For each bank take its entire raw interval of $K$ labels in
\eqref{eq:translated-grid-label}.  A raw label from one bank and one
from a different bank coincide with probability exactly $1/p$:
their difference is a nontrivial affine equation in $u_0,v_0$, with
coefficient $\theta-\phi\ne0$ at $u_0$.  If $C$ counts unordered
cross-bank raw-label collisions, then
\[
 \mathbb E C\le\binom L2\frac{K^2}{p}=o(LHM),
 \qquad \frac{LHM}{p}=O((\log H)^{-3}).
\]
Markov's inequality gives $C\le LHM/64$ with probability $1-o(1)$.
Hence some translation satisfies this and $X\ge M^2/4$ simultaneously.
Of the at least $LHM/16$ rich bank-label pairs, discard every pair
involved in any cross-bank raw-label collision.  At least $LHM/32$
remain, each with a unique bank owner even among non-rich fibers.

The union $Z$ of raw label intervals has size at most $LK=o(p)$.
Choose
\[
 c_0\notin Z\cup(Z-1)\cup\{0,-1\}.
\]
Then labels zero and one have no bank matches on the grid.  This
shift does not affect fiber sizes or collision counts.

\paragraph{Neutral padding and exclusion of other quadratics.}
Set $T=A+d$ and
\[
 n=\lfloor T^2/2\rfloor+1.
\]
For large $H$, append $n-N_0-X>0$ neutral coordinates, setting
$g(x)=0$ and $f(x)=x^3$.  Avoid the old domain, zero, and the at most
$3L$ roots of $x^3-P_\theta(x)$ over all banks.  There are enough
coordinates since $p\gg n=\Theta(L^2)$.  Banks gain no neutral matches.

Fix a label $\lambda$.  For any nonconstant quadratic $Q$, the grid
equation $Q(x)=(c_0+\lambda)/U$ has at most two roots for each of the
$M$ possible values of $U$, hence at most $2M$ grid matches.
A nonzero constant $Q$ has matches in at most one such column and
hence at most $M$ matches.  The sole exception is $Q=0$ at
$\lambda=-c_0$, which we exclude.  Every other nonbank quadratic has
at most
\[
 L+2M+3<A
\]
matches on the entire domain, the last three accounting for the neutral
coordinates.  Consequently every retained rich label except possibly
$-c_0$ has a singleton list at threshold $T$: its unique bank owner
has at least $A+d$ matches, other banks have only their $A$ core
matches, and all nonbank words have fewer than $A$ matches.
There are at least
\[
 B=\left\lfloor\frac{LHM}{32}\right\rfloor-1
\]
such labels.  Labels zero and one have exact maximum agreement $A$.
At every label all $L$ bank words meet the relaxed threshold $A$;
the only possible additional word is zero at $-c_0$.  Since
$X\ge M^2/4>A$ eventually, the relaxed list has size exactly $L+1$
at that exceptional label and exactly $L$ elsewhere.

For common agreement, a nonzero quadratic explaining $g$ has at
most two matches on the zero-direction coordinates and at most
$2M$ on the grid.  Its total is below $A$.  The zero explanation
restricts common matches to the core and neutral coordinates; there
the maximum agreement of $f$ is exactly $A$, attained by every bank
word.  Thus $\operatorname{CA}_3(f,g)=A$.
Finally set $F=f$, $G=f+g$.  Both individual agreements are $A$,
and the invertible change of the two source coordinates preserves
common agreement.  Their affine mixtures are exactly $f+\lambda g$.

\paragraph{Parameter comparison.}
We have $M=\Theta(L)$, $d=\Theta(L/H)=o(L)$, and
\[
 n\sim2L^2=\Theta\!\left(\frac{H^3}{(\log H)^3}\right).
\]
Therefore $H=\Theta(n^{1/3}\log n)$, giving all displayed scales.
The equality $n=\lfloor T^2/2\rfloor+1$ implies $T^2<2n$ exactly.
Moreover $A/\sqrt n\to\sqrt2$, whereas the dimension-three
first-order agreement bound is
\[
 n a_1(3/n)\le\sqrt{3n/2}+(3n/8)^{1/4}.
\]
Hence $n a_1(3/n)<A$ eventually, proving the location of the example.
\end{proof}

The grid uses the incidence budget efficiently: with
$\varepsilon=d/n$ and relaxed line-local list bound $L+1$, its
singleton count is $\Omega((L+1)/\varepsilon)$.  This is a statement
about the displayed coupled parameters, not a universal decoding bound
or a code-wide list-size estimate.  The characteristic grows much faster
than the message degree, while the code rate tends to zero.
Conversely, apart from the single zero-polynomial exception, every
successful label consumes at least $d$ of the $LX\le LM^2$ fresh
bank-label incidences.  The total exceptional fraction is therefore
$O((\log H)^{-3})$, consistent with the theorem's scope.

\begin{corollary}[A range of correlated gaps]
\label{cor:translated-grid-gap-range}
There is an absolute constant $c>0$ such that the same construction
works whenever $H\to\infty$ and integers $L$ satisfy
\[
 L/H\longrightarrow\infty,\qquad
 L\le c(H/\log H)^{3/2}.
\]
With $M=\lfloor L/4\rfloor$, $d=\lfloor M/(64H)\rfloor$, and the
same exact choices of $A,T,n$, it gives
\[
 n\sim2L^2,\qquad
 B=\Omega(n^{3/2}/d),\qquad
 p=\Theta(n^2/d^4),
\]
with the same source, common-agreement, and singleton-list conclusions.
\end{corollary}
\begin{proof}
Retain any $L$ edges of the preceding four-cycle-free graph.
Its Sidon property and all root counts persist.  The retention and
cross-label estimates still hold, because $L/H\to\infty$ and
$LHM/p=O(L^2/H^3)=O((\log H)^{-3})$.
Now $d=\Theta(L/H)$, $B=\Omega(HL^2)$, and $p=\Theta(H^4)$;
substituting $n\sim2L^2$ proves the formulas.
\end{proof}
